\documentclass[
  12pt,
  a4paper,
  oneside,
  brazil
]{abntex2}

% -------------------------------------------------------
% Pacotes
% -------------------------------------------------------
\usepackage[T1]{fontenc}
\usepackage[utf8]{inputenc}
\usepackage{lmodern}
\usepackage{microtype}
\usepackage{indentfirst}
\usepackage{hyperref}
\usepackage{abntex2cite}

% -------------------------------------------------------
% Metadados
% -------------------------------------------------------
\titulo{Perfil do usuário do tema ``Segurança Virtual / Identificação de Sites Falsos''}
\autor{} % preencher
\instituicao{} % preencher
\data{\today}

% -------------------------------------------------------
\begin{document}

\pretextual
\maketitle

\textual

% =======================================================
\section{Características demográficas}
% =======================================================

O público-alvo desse tema abrange principalmente adolescentes e jovens (10--29 anos), bem como
pessoas a partir de 40--60+ anos, com maior vulnerabilidade entre idosos que migraram
recentemente para o uso digital \cite{portalenvelhecimento2025}. A localização típica é o
Brasil, em cidades e regiões urbanas e algumas rurais com acesso à internet via celular e banda
larga \cite{ibge_tic2023}. Em relação à escolaridade, observa-se maior dificuldade entre
usuários com ensino fundamental incompleto ou sem instrução, embora jovens com ensino médio
também apresentem baixa alfabetização digital \cite{ibge_pnad_educ2023, educafcc2022}.
Estima-se que cerca de 88\% das pessoas com 10 anos ou mais já usem internet no Brasil, com
crescimento acelerado principalmente entre idosos \cite{ibge_tic2023}.

% =======================================================
\section{Comportamentos e hábitos relacionados ao tema}
% =======================================================

O usuário costuma acessar, pelo smartphone, bancos, redes sociais, WhatsApp, lojas virtuais e
serviços de governo, como parte de sua rotina diária \cite{bcb_rig2024, ibge_tic2023}. Muitas
vezes, navega sem verificar se o site é verdadeiro ou falso, assumindo que, se o aplicativo ou
a página parecem iguais aos legítimos, então são seguros \cite{certbr2024}. Com frequência,
não utiliza senhas fortes, reutiliza senhas em diferentes serviços e clica em links recebidos
por mensagem instantânea sem confirmar a origem \cite{bcb_rig2024}.

% =======================================================
\section{Necessidades e dificuldades}
% =======================================================

A necessidade central desse público é ter uma forma simples de saber se o site ou o link é
seguro antes de informar dados pessoais ou financeiros, como senha, CPF ou número de cartão
\cite{abecs2023, certbr2024}. A principal dificuldade está em não reconhecer sinais de
\textit{phishing}, como URL estranha, domínio diferente ou com letras trocadas, bem como não
perceber quando o cadeado de HTTPS não é suficiente para garantir segurança \cite{certbr2024}.
Idosos e usuários com baixa escolaridade costumam ter baixa alfabetização digital, apesar de já
usarem bastante internet, o que aumenta a chance de cair em golpes
\cite{ibge_pnad_educ2023, educafcc2022}.

% =======================================================
\section{Objetivos e motivações}
% =======================================================

Os objetivos principais do usuário são conseguir realizar transferências, compras e pagamentos
sem risco de perda financeira e proteger seus dados pessoais (CPF, RG, fotos, renda) de
vazamento e uso indevido \cite{bcb_rig2024, abecs2023}. As motivações incluem a comodidade de
acessar tudo pelo celular e o medo de ser enganado, diante da grande divulgação de golpes de
\textit{phishing} na mídia e nas redes sociais \cite{abecs2023, certbr2024}.

% =======================================================
\section{Contexto de uso}
% =======================================================

O uso ocorre, em geral, em casa, usando smartphone com Wi-Fi ou dados móveis, e, em alguns
casos, em \textit{lan houses} ou equipamentos compartilhados \cite{ibge_tic2023}. Com
frequência, o acesso acontece em momentos de urgência, como quando há boleto vencido, oferta
relâmpago ou mensagem alegando ``problema no banco'' \cite{abecs2023, bcb_rig2024}. Normalmente,
o usuário clica em links recebidos por WhatsApp, SMS, e-mail ou redes sociais, muitas vezes sem
conferir a origem ou analisar o domínio do site aberto \cite{certbr2024}.

% =======================================================
\section{Principais dores e expectativas}
% =======================================================

Entre as dores principais estão ser enganado por site falso de banco, loja ou serviço
governamental, resultando em roubo de dinheiro ou clonagem de cartão, além de não saber o que
fazer após perceber o golpe, como bloquear o cartão ou comunicar o banco adequadamente
\cite{abecs2023, bcb_rig2024}. Também é comum a sensação de vergonha ou medo de falar sobre o
erro, o que faz com que muitos casos sequer sejam denunciados \cite{abecs2023}. Em relação a
uma possível solução, o usuário espera uma verificação de site simples e intuitiva, com poucos
cliques, acompanhada de alerta claro e imediato (``esse site é falso'' ou ``suspeito'') ao
tentar digitar a senha \cite{certbr2024}. Além disso, espera conteúdo educativo curto e visual,
com ícones e exemplos, para poder reconhecer sites falsos nas próximas vezes
\cite{educafcc2022}.

% =======================================================
\postextual
\bibliography{referencias}

\end{document}
